

O crescimento exponencial da produção científica representa um desafio crescente para a descoberta e curadoria de conhecimento. Repositórios como o arXiv \cite{arxivorg_submitters_arxiv_nodate} acumulam milhões de artigos, com novos trabalhos submetidos diariamente em áreas como computação, física e matemática. Nesse contexto, sistemas capazes de gerar automaticamente metadados, como títulos e palavras-chave, a partir do conteúdo textual dos artigos têm aplicação direta em indexação, busca semântica e recomendação acadêmica.

Com o advento dos Grandes Modelos de Linguagem (LLMs), tornou-se possível abordar essa tarefa com qualidade sem precedentes. Entretanto, o emprego direto de LLMs pré-treinados sem especialização produz resultados inconsistentes: conforme observado neste estudo, o modelo sem ajuste fino gerou títulos excessivamente longos e com artefatos do \textit{prompt} de instrução, evidenciando a necessidade de especialização para a tarefa. O principal obstáculo ao ajuste fino em contextos com restrições computacionais é o custo proibitivo do treinamento completo, que exige o armazenamento de gradientes e estados do otimizador para todos os parâmetros do modelo. Métodos de ajuste fino com eficiência paramétrica (PEFT)~\cite{ding_parameter-efficient_2023}, como LoRA~\cite{hu_lora_2022} e Q-LoRA~\cite{dettmers_qlora_2023}, surgem como alternativas que reduzem drasticamente esse custo. Contudo, a comparação sistemática dessas estratégias em tarefas aplicadas de geração de texto com modelos de pequeno porte permanece pouco explorada na literatura.

Este trabalho investiga esse problema ao comparar quatro estratégias de adaptação (Baseline zero-shot, ajuste da última camada (\textit{Linear Probing}), Q-LoRA e \textit{Full Fine-Tuning}) aplicadas ao modelo Qwen2,5-0,5B-Instruct para a geração automática de títulos a partir de resumos do arXiv. Os resultados demonstram que o Q-LoRA alcança desempenho estatisticamente equivalente ao treinamento completo ($p = 0,109$), validando sua viabilidade como alternativa eficiente para pesquisadores e instituições com infraestrutura limitada.

A tarefa se assemelha à sumarização automática de textos, na qual um LLM processa um texto de entrada para retornar uma versão reduzida que preserva o sentido do conteúdo original. Entretanto, o objetivo difere ligeiramente: ao invés de produzir resumos, trata-se da geração de títulos a partir de resumos de artigos científicos. Esta tarefa apresenta desafios intrínsecos: é comum que autores optem por títulos ``chamativos'' ou estilizados que não podem ser diretamente inferidos a partir do resumo. Quando tais títulos são utilizados como \textit{ground truth} no treinamento, o modelo enfrenta dificuldades de convergência, pois a correlação semântica entre a entrada e a saída esperada é fraca.

Outro desafio crítico reside na métrica de avaliação. A comparação literal entre o título gerado e o original é insuficiente para capturar a qualidade da predição. Considere dois cenários de divergência: o título predito é uma permutação das palavras do título original; ou o título predito utiliza vocabulário distinto, mas com o mesmo significado semântico do original.

Para a aplicação proposta, o segundo caso é preferível, pois demonstra a capacidade do modelo de abstrair e sintetizar o conteúdo, em vez de apenas memorizar sequências de tokens. Torna-se necessário, portanto, empregar uma métrica que avalie a semelhança semântica. Uma abordagem amplamente adotada é o Sentence-BERT~\cite{reimers_sentence-bert_2019,devlin_bert_2019}, um modelo baseado em BERT treinado especificamente para medir similaridade semântica entre sentenças, gerando valores próximos de 0 para frases sem relação semântica e próximos de 1 para frases semanticamente equivalentes.

O restante deste artigo está organizado da seguinte forma. A Seção~\ref{sec:related} apresenta os trabalhos relacionados. A Seção~\ref{sec:metodologia} descreve a metodologia adotada, incluindo a construção do \textit{dataset} e as técnicas de treinamento avaliadas. A Seção~\ref{sec:experimento} detalha o protocolo experimental e as métricas de avaliação. A Seção~\ref{sec:resultados} apresenta e discute os resultados obtidos. Por fim, a Seção~\ref{sec:conclusao} sintetiza as conclusões e aponta direções para trabalhos futuros.